\documentclass[12pt]{article}
\usepackage[utf8]{inputenc}
\usepackage[T1]{fontenc}
\usepackage{hyperref}
\title{Pambia Language Resources for K-12 Curriculum Development}
\author{Compiled by Claude Sonnet 5.5 (claude-sonnet-5-5) for Andr\'as Kornai}
\date{}
\begin{document}
\maketitle

\section{Introduction}
Pambia (ISO 639-3 \texttt{pmb}, Glottocode pamb1243) is an extinct Zande language of northeastern Democratic Republic of the Congo. The only speaker count found is 21,000, from an older Ethnologue edition (1982). The sources found here give no extinction date. One database (ASJP) lists it as extinct since 2024. The search found no teaching, revitalization or documentation materials beyond the basic reference entries below.

\section{Primary Educational Resources}
None were found.

\section{Linguistic and Reference Materials}
\begin{itemize}
\item Wikipedia article on Pambia: \url{https://en.wikipedia.org/wiki/Pambia_language}
\item Wikipedia article on the Zande languages: \url{https://en.wikipedia.org/wiki/Zande_languages}
\item ASJP entry for Pambia: \url{https://asjp.clld.org/languages/PAMBIA}
\item Grambank entry for Pambia: \url{https://grambank.clld.org/languages/pamb1243}
\end{itemize}

\section{Religious and Community Resources}
None were found.

\section{Digital and Audio Resources}
None were found.

\section{Curriculum Development Context}
A heritage curriculum for Pambia would depend on the closely related Zande (Azande) language and its materials. These were not searched for here. The limited Pambia data (ASJP and Grambank entries) would have to be checked against fuller sources before use.
\end{document}